% ============================================================
% 04-experiments.tex — V13 Experimental Setup (condensed)
% ============================================================
\section{Experimental Setup}
\label{sec:exp_setup}

%------------------------------------------------
\subsection{Dataset and Annotation}
\label{sec:dataset}

We evaluate \system{} on 200 administrative queries constructed through a document-grounded protocol over official CTU statutes. The 100-query development split supports index calibration and prompt engineering. The 100-query held-out set follows a predefined distribution: 40 direct, 20 within-domain multi-hop, 20 cross-domain, 10 entity-comparison, 5 temporal, and 5 adversarial queries (51 formal, 49 colloquial; 56 single-source, 44 multi-source). It contains 19 Academic, 28 Financial, 17 Scholarship, and 16 General single-domain queries; the remaining 20 combine evidence across domains. Candidate queries and ground-truth evidence were extracted from official statutes: one annotator established reference answers, required facts, and gold sources, while a second annotator independently verified all annotations; disagreements were adjudicated against official regulatory text. An audit found no cross-split duplicates (max 5-gram Jaccard = 0.1875). Each query is annotated along nine dimensions, including domain, evidence type, and tool requirements, to support fine-grained failure analysis.

%------------------------------------------------
\subsection{Baselines and Configurations}
\label{sec:baselines}

\textbf{Fairness principle and execution contract.} For Scenario~1, all architecture configurations share the same LLM (Gemini 2.5 Flash Lite, $T{=}0.0$), knowledge representations, retrieval stack (BM25 + dense + RRF + \texttt{bge-reranker-v2-m3}, $k{=}7$), and deterministic calculator access. Differences are strictly restricted to \emph{orchestration architecture}. Scenario~2 intentionally changes one component at a time. To isolate orchestration differences from retrieval variance in Scenario~1, all three architecture variants evaluate over the identical shared full-stack evidence retrieval output, yielding matched pre-generation Source Recall (78.6\%) and Source AP (70.8\%). The shared evidence context excludes outputs of tools whose invocation is evaluated; tool-derived results are available only after successful tool execution. Performance divergences in final answers are therefore attributable purely to tool privilege scoping, prompt specialization, and agent routing. Multi-agent variants use intent-routed scoped registries with identical execution bounds: maximum 3 tool calls, matching context limits, 2{,}048 output-token budget, and 30-second timeouts. Token accounting for \system{} comprehensively aggregates all pipeline invocations (supervisor routing, contract repair, worker generation, and retries), ensuring rigorous parity with the single-agent baseline.

\paragraph{Scenario 1 --- Architecture Comparison (RQ1):}
Three configurations are evaluated on the held-out set ($N{=}100$ queries, 300 runs under greedy decoding $T{=}0.0$):
\textbf{S1-A:} Unified Single Agent with all 11 tools, no supervisor or routing;
\textbf{S1-B:} Routed Generic Worker---supervisor routes to domain workers with scoped tools but a single generic prompt;
\textbf{S1-C:} Full CTU-Chat with supervisor, contract repair, specialist prompts, scoped registries, and deterministic computation.

\paragraph{Scenario 2 --- Ablation (RQ2):}
Five ablations each remove one component from Full CTU-Chat ($N{=}100$ queries per ablation, $T{=}0.0$):
\textbf{S2-A1:} No Specialist Policy (generic worker prompt);
\textbf{S2-A2:} No Tool Isolation (all 11 tools visible);
\textbf{S2-A3:} Uniform Text-RAG Baseline (graph data serialized to text, traversal disabled);
\textbf{S2-A4:} No Deterministic Calculation (LLM generates arithmetic directly);
\textbf{S2-A5:} No Route Repair (raw supervisor output used directly).

\paragraph{Scenario 3 --- Tool Ambiguity (RQ3):}
Fifty fixed cases are evaluated while 40 semantic-neighbor tools expand registries to 11, 31, and 51 tools. Four configurations are compared: Full-Registry Single Agent, Global Top-10 Single Agent, Routed Generic, and Routed Specialist. Each case uses 3 counterbalanced tool-order blocks, yielding $50{\times}4{\times}3{\times}3{=}1{,}800$ decisions across counterbalanced blocks.

%------------------------------------------------
\subsection{Evaluation Metrics}
\label{sec:metrics}

Metrics span four groups. \emph{Final-answer quality:} Required-Fact Coverage (continuous proportion of reference facts verified in generated response), Answer Correctness, Answer Relevancy, and Faithfulness (Ragas~\cite{es2023ragas}). \emph{Evidence quality:} Source Recall, Source AP, Context Recall, Context Precision. \emph{Execution reliability:} Tool Invocations per Query, Tool Selection Accuracy, Argument EM/Contract Match, Result Accuracy, and route repair rate. \emph{Operational cost:} input/output tokens, LLM call counts, and wall-clock latency. In comparative evaluation, E2E Success Rate serves as the primary endpoint. Key secondary endpoints are Required-Fact Coverage, tool invocations, token cost, and latency. Source Recall and Source AP serve as controlled retrieval diagnostics, expected identical by design across Scenario~1 configurations.

\textbf{E2E Success definition.} For retrieval queries, E2E Success is evaluated under the criterion of required facts ${\ge}50\%$, absence of unsupported claims, and retrieval of verified sources, complemented by continuous Fact Coverage. For tool queries: correct route, tool, arguments, result, and answer; for the Unified Single Agent baseline, routing correctness is evaluated as N/A by construction as it does not employ a router. For adversarial queries: absence of unnecessary tool calls and appropriate clarification or refusal. For the arithmetic subset ($N{=}9$), numeric exact match against official bursar schedules is independently audited.

%------------------------------------------------
\subsection{Statistical Protocol}
\label{sec:statistics}

Evaluations use greedy decoding ($T{=}0.0$) to minimize sampling variance. Paired differences across configurations are computed at the query level ($N{=}100$); 95\% CIs are estimated using 10,000 bootstrap resamples. The primary confirmatory hypothesis is CTU-Chat versus Unified Single Agent on E2E Success. Generic-worker comparisons, continuous Fact Coverage, token cost, and latency are evaluated as secondary outcomes with reported effect sizes and 95\% bootstrap CIs. Scenario~3 degradation is computed as $\Delta = \text{score}_{51} - \text{score}_{11}$ across counterbalanced tool-order blocks.

%------------------------------------------------
\subsection{Infrastructure}
\label{sec:infrastructure}

Gemini 2.5 Flash Lite via Vertex AI ($T{=}0.0$). Neural reranking on a GTX 1650 (evaluation hardware). Neo4j~5.20 stores relational curricula; PostgreSQL stores parent chunks (2,800-char, 100 overlap); Qdrant indexes child chunks (400-char, 50 overlap) with a 768-dimensional \texttt{vietnamese-bi-encoder}. System uses \texttt{bge-reranker-v2-m3}, LangGraph StateGraph, and Redis~7. Latency excludes the first warm-up run. Prompt templates, tool contracts, and graph schema are archived.
